\documentclass[12pt]{article}
\usepackage[T1]{fontenc}
\usepackage[utf8]{inputenc}
\usepackage{lmodern}
\usepackage{textcomp}
\usepackage{enumitem}
\usepackage{geometry}
\geometry{margin=1in}
\begin{document}
\begin{center}
Answer Key - Computer Science - Graduate Level Assessment 6 \
\small (Answers are ordered exactly as in the question paper.)
\end{center}

\section*{Multiple Choice}
\begin{enumerate}[label=\arabic*.]
\item \textbf{Answer:} B
\\ \textit{Options:} A) Physical layer; B) Application layer; C) Data link layer; D) Network layer
\\ \textit{Note:} A proxy firewall operates at the application layer because it intercepts and analyzes the traffic between a client and server, allowing it to filter content based on application-specific protocols and data.
\item \textbf{Answer:} B
\\ \textit{Options:} A) Confidentiality; B) Availability; C) Correctness; D) Integrity
\\ \textit{Note:} Availability is not considered a classic security property like confidentiality and integrity, as it focuses on ensuring access to information rather than the protection of that information itself.
\item \textbf{Answer:} C
\\ \textit{Options:} A) Eavesdropping; B) Cross-site scripting; C) Authentication; D) SQL Injection
\\ \textit{Note:} Authentication is a fundamental security process used to verify the identity of a user or system, rather than a security exploit, which typically involves malicious actions aimed at compromising security.
\item \textbf{Answer:} A
\\ \textit{Options:} A) 1/K; B) K-1/K; C) log\_10 (1/N); D) N-1/m
\\ \textit{Note:} In a complete K-ary tree of depth N, the total number of nodes is given by ($K^(N$+1) - 1) / (K - 1), while the number of nonterminal nodes, which are all internal nodes, is $K^N$ - 1, leading to the ratio of nonterminal nodes to total nodes being approximately 1/K as N becomes large.
\item \textbf{Answer:} B
\\ \textit{Options:} A) Memory leakage; B) Buffer-overrun; C) Less processing power; D) Inefficient programming
\\ \textit{Note:} A buffer-overrun occurs when a program writes more data to a buffer than it can hold, leading to potential overwriting of adjacent memory, which can be exploited by attackers to gain unauthorized access or disrupt system functionality.
\end{enumerate}

\section*{True / False}
\begin{enumerate}[label=\arabic*.]
\setcounter{enumi}{5}
\item \textbf{Answer:} True
\\ \textit{Options:} A) True; B) False
\\ \textit{Note:} Wireless security encompasses protocols, technologies, and practices designed to safeguard data integrity and prevent unauthorized access to devices connected through wireless networks.
\item \textbf{Answer:} False
\\ \textit{Options:} A) True; B) False
\\ \textit{Note:} Church's thesis, which posits that any effectively calculable function can be computed by a Turing machine, remains unprovable as it is a foundational assumption in computer science rather than a theorem that can be definitively proven or disproven with advancements in computational theory.
\item \textbf{Answer:} False
\\ \textit{Options:} A) True; B) False
\\ \textit{Note:} WPA 2 is not a recent protocol; it was established in 2004 as an enhancement to WPA and is widely used for securing wireless networks, rather than being a new feature that emerged in recent years.
\item \textbf{Answer:} True
\\ \textit{Options:} A) True; B) False
\\ \textit{Note:} The set of Boolean operators \{AND, OR\} is not complete because it lacks the ability to express negation, which is essential for representing all possible Boolean expressions.
\item \textbf{Answer:} False
\\ \textit{Options:} A) True; B) False
\\ \textit{Note:} A front-door typically refers to a legitimate access point in a system, but it may not always be fully traceable or compliant with security protocols, as vulnerabilities can exist even in authorized access methods.
\end{enumerate}

\section*{Long Form Response}
\begin{enumerate}[label=\arabic*.]
\setcounter{enumi}{10}
\item \textbf{Answer:} def alternate\_elements(list 1): | return result
\\ \textit{Note:} correct because it outlines the essential function structure needed to generate a new list containing alternate elements from the provided list, indicating an understanding of list manipulation and function definition in programming.
\item \textbf{Answer:} def sort\_mixed\_list(mixed\_list): | return int\_part + str\_part
\\ \textit{Note:} The answer correctly outlines a function that separates integers and strings from a mixed list and returns them concatenated, demonstrating an understanding of data types and list manipulation in programming.
\end{enumerate}

\vfill
\noindent\textit{End of Answer Key}
\end{document}